\chapter{Dataset and Experimental Setup}
\label{ch:setup}

%-------------------------------------------------------

The setup below applies to all five tasks. The code is split into five
Python files.
\texttt{dataloader.py} downloads MNIST and builds the data split,
\texttt{model.py} defines the CNN, \texttt{train.py} contains the
training and evaluation loops, and \texttt{plots.py} draws every figure.
\texttt{main.py} runs all five tasks in order and writes a log of the
program output to \texttt{results/output\_log.txt}. I started from the
``CNN for MNIST'' example in the lecture notes \citep{abidl08}, which is
based on the PyTorch tutorials \citep{pytorchtutorial}. I adapted it to
the $28 \times 28$ MNIST images and added the validation split, the
metrics and the experiment.


%%%%%%%%%%%%%%%%%%%%%%%%%%%%%%%%%%%%%%%%%%%%%%%%%%%%%
\section{Dataset and Data Split}
\label{sec:data}
%%%%%%%%%%%%%%%%%%%%%%%%%%%%%%%%%%%%%%%%%%%%%%%%%%%%%

MNIST \citep{lecun1998,mnistkaggle} contains $28 \times 28$ grayscale
images of the handwritten digits 0--9: 60,000 for training and 10,000 for
testing. I load it with \texttt{torchvision.datasets.MNIST}, which
provides the same files as the Kaggle version. The 60,000 training images
are split once into 48,000 for training (80\%) and 12,000 for validation
(20\%) using a random permutation drawn from a generator seeded with 42
(Listing \ref{lst:split}). Every experiment calls the same function with
the same seed, so all of them use the exact same split. The 10,000 test
images are not used for anything until Task 5.

\lstinputlisting[language=Python, firstline=18, lastline=33,
    firstnumber=auto,
    caption={Fixed 80/20 split and normalization statistics (\texttt{dataloader.py}).},
    label={lst:split}]{../dataloader.py}

\textbf{Preprocessing.} Each image goes through two transforms:
\texttt{ToTensor}, which turns the 0--255 pixel values into a
$1 \times 28 \times 28$ float tensor in $[0, 1]$, and \texttt{Normalize},
which subtracts the mean and divides by the standard deviation. I
computed the mean (0.1307) and standard deviation (0.3082) from the
48,000 training images only, so the validation and test images do not
influence the preprocessing. No data augmentation is used. Normalized
inputs are centered around zero, which keeps the first layer's outputs
in a reasonable range at the start of training.


%%%%%%%%%%%%%%%%%%%%%%%%%%%%%%%%%%%%%%%%%%%%%%%%%%%%%
\section{Random Seed and Identical Starting Conditions}
\label{sec:seed}
%%%%%%%%%%%%%%%%%%%%%%%%%%%%%%%%%%%%%%%%%%%%%%%%%%%%%

All experiments use the seed 42. \texttt{set\_seed} (Listing
\ref{lst:seed}) seeds Python, NumPy and PyTorch and turns on PyTorch's
deterministic algorithms. It is called before the network is created, so
every configuration starts from the same initial weights. The training
loader shuffles with its own generator, also seeded with 42, so all runs
see the batches in the same order. Data loading uses
\texttt{num\_workers=0}, which keeps everything in one process.

\lstinputlisting[language=Python, firstline=24, lastline=29,
    firstnumber=auto,
    caption={Seeding every random number generator (\texttt{train.py}).},
    label={lst:seed}]{../train.py}
